\documentclass{article}

\usepackage{arxiv}

\usepackage[utf8]{inputenc} % allow utf-8 input
\usepackage[T1]{fontenc}    % use 8-bit T1 fonts
\usepackage{hyperref}       % hyperlinks
\usepackage{url}            % simple URL typesetting
\usepackage{booktabs}       % professional-quality tables
\usepackage{amsfonts}       % blackboard math symbols
\usepackage{amsmath}
\usepackage{microtype}      % microtypography

\title{A CRISP-DM Description for Anomaly Detection in Industrial and Satellite Telemetry}

\author{
 Alessandro Bevilacqua \\
  Department of Computer Science \\
  University of Bari Aldo Moro \\
  Bari, Italy \\
  \texttt{a.bevilacqua11@studenti.uniba.it} \\
}

\begin{document}
\maketitle

This document reports the \emph{process}: the decisions made in each CRISP-DM phase, 
their rationale, and the artifacts passed to subsequent stages. 
Crucially, every decision described below is encoded in a DVC pipeline, ensuring that each step can be verified directly against its resulting artifact.

\section{Business Understanding}

Industrial facilities and spacecraft systems generate continuous multivariate telemetry data, where operational faults are rare but incur extremely high costs.

\paragraph{Business objectives.} 
The primary operational objectives reflect real-world monitoring requirements. 
First, the detection system must identify anomalies in real time, 
enabling preventive intervention before catastrophic failure occurs rather than post-hoc analysis. 
Second, it must maintain a low false alarm rate to minimize operator fatigue, given that operational attention is a strictly 
resource in control rooms. 
Third, it must maintain a low missed anomaly rate, since an anomaly can cause significant damage to downstream equipment.


\paragraph{Business success criteria.} 
In the absence of domain experts for direct evaluation, operational requirements are translated into quantitative metrics.
The first metric measures the proportion of annotated failure episodes identified by the model. 
The second metric assesses precision by measuring the ratio of true anomalies among all raised alarms. 
Since both metrics depend on the detection threshold, they are reported alongside threshold-independent performance measures to provide an unbiased comparison.

\paragraph{Data mining goals.} 
The operational objectives translate into two main technical goals. 
The first is to construct an \emph{inductive} model of normality capable of scoring incoming real-time observations without requiring re-training. 
The second and central research goal is to empirically assess whether an observation's structural neighbourhood derived from clustering provides informative 
context that an isolated algorithm cannot capture. 
The objective is to provide a rigorous, defensible answer to this question,
where a negative finding represents a valid scientific result.

This research question dictates the experimental setup. The proposed graph-based detector (DENCAD) is evaluated against four controlled baselines:
\begin{enumerate}
    \item Plain autoencoder: The identical network architecture with context removed (reducing adjacency to the identity matrix);
    \item Graph autoencoder: The same network evaluated over an $L$-hop neighbourhood without clustering;
    \item K-Means: A centroid-based baseline;
    \item DENCAST distance: A density-based clustering baseline using the nearest clustered object distance.
\end{enumerate}

\paragraph{Project plan.} 
The execution follows these structured phases: 
analyzing raw data sources and labels; 
constructing training, validation, and test splits across labelling scenarios using a temporal approach;
fitting models on the training set; 
calibrating hyperparameters on the validation set; 
refitting all five detectors using the combined train-validation data; 
and evaluating final performance on an isolated test set. 
Every phase is fully versioned and tracked via a DVC pipeline to ensure reproducibility.

\section{Data understanding}

Three public datasets were selected to ensure that the conclusions do not rely on a single operational regime:
two industrial testbeds with annotated physical attacks and one spacecraft telemetry archive containing real operational faults. 
Each dataset was profiled prior to pipeline development to address three fundamental questions: 
feature composition, ground-truth label structure, and the exact definition of normal versus anomalous behavior.

\paragraph{SWaT: The water treatment testbed.} A single tabular recording at one reading per
second, labelled normal or under attack over a campaign of staged intrusions. Which
channels hold a discrete state rather than a measurement is declared by the instrument
naming convention, so the distinction is read from the documentation.
The profiling's substantive finding is
that some actuators are never recorded at all during the earliest part of the recording,
which is also the part that will become the training period: a property invisible in the
dataset as a whole and visible only once the split is contemplated.

\paragraph{HAI: The hardware-in-the-loop dataset.} Several recordings and a separate label file
to be joined on their timestamps, again at one reading per second. The discrete channels
are taken from the table in the dataset's own technical documentation.
Profiling then showed that most of the declared discrete channels never
vary at all, and that several more are one selector read under different names, so only a
few of them survive to carry any information.

\paragraph{ESA: The spacecraft archive.} The only one of the three that is not a table to begin
with. Each channel is a separate serialised series on its own irregular index, sampled at
a rate that differs from channel to channel, and the labels are not a column but a set of
annotated intervals, each attached to a particular channel and carrying a category.
The categories separate anomalies from \emph{rare events}: unusual but
legitimate behaviour, annotated precisely so that it would not be counted as a fault. 
The profiling established a heavy duplication among the channels, with many pairs that are plainly the same quantity reported twice.

\paragraph{ESA: Selecting target mission.} The ESA telemetry archive encompasses multiple missions, 
requiring a deliberate selection strategy.
Given the strict chronological nature of the train-validation-test split, 
a mission is viable only if a sufficient number of annotated anomaly episodes occur \emph{after} the initial training window. 
Missions where anomalies are clustered within the early operational phases fail this criterion regardless of total data volume, 
as the critical parameter is not total telemetry count but the availability of distinct anomalous episodes across validation and test partitions. 
Consequently, Mission 1 was selected for this study and all the others were discarded.

\section{Data preparation}

Each dataset is processed through a dedicated preprocessor.
Two abstract base classes define the general preprocessing pipeline and common operations, 
while specific subclasses extend them to implement dataset-customized transformations.

\paragraph{Assembly.} Tabular datasets undergo deduplication,
chronological sorting, and timestamp-based concatenation where recordings span multiple files.
Conversely, the ESA telemetry archive requires temporal resampling: 
all channels are mapped onto a uniform time grid via zero-order hold interpolation—holding each channel's latest value until a new observation is logged—while 
continuous annotation intervals are converted into discrete per-row labels. 
The sampling grid interval is thus treated as an explicit modeling choice.

\paragraph{Cleaning: which columns survive.} At assembly, rows null in every feature are
dropped, as are rows identical to an earlier one including the timestamp. The processing
step then applies three column filters in a fixed order, each decided on the training rows.
First, any column missing more than a set share of them is dropped, which is what removes
channels never recorded during the training period. Second, any column constant across the
whole training split is discarded. Third, a redundancy filter removes one of every
pair that is the same quantity reported twice, judged by linear correlation for numeric
pairs and by Cram\'er's $V$ for categorical ones, both against a threshold just below
unity. 
Which of the two copies survives is decided rather than left to column order:
numeric columns are considered in order of decreasing number of distinct values and
categorical ones in order of decreasing normalised entropy.

\paragraph{Imputation.} The fill values are fitted only after those
drops, on the training rows that remain: the median of each numeric column and the mode of
each declared categorical one.

\paragraph{Data splitting.} Data splitting is performed strictly chronologically across all three datasets,
ensuring that the validation and test sets temporally succeed the training period and that no sliding instances overlap across partitions.

\paragraph{Windowing.} Readings are aggregated into non-overlapping windows of a fixed
number of readings, chosen per dataset, and the window rather than the reading is the unit
that is scored and labelled: a window is positive if any of its readings is. Numeric
channels contribute their mean, standard deviation, minimum and maximum, discrete ones
their mode, entropy, number of distinct values and share of the dominant value.
Additionally, windows containing fewer than half of the expected readings are discarded.


\paragraph{Scaling.} Window features are standardised with the mean and deviation of the
training windows.

\paragraph{Training composition and labeling regimes.} 
The training splits for the SWaT and HAI datasets contain no anomalous instances, 
providing a purely nominal baseline. 
Conversely, the training partition of the ESA telemetry archive contains residual anomaly occurrences. 
To evaluate the impact of training contamination, two experimental setups are established: 
one in which these anomalous instances are removed prior to training, and another in which they are retained.

Regarding target annotations, 
SWaT and HAI enforce a strict binary scheme (normal versus anomalous). 
The ESA dataset, by contrast, differentiates between operational anomalies and rare system events. 
To assess label sensitivity, both interpretations are evaluated: one treating rare events as normal operating conditions, 
and the other categorizing them as anomalous episodes.
Combined with the two training contamination choices, these target definitions formulate four distinct pipeline variants
for the ESA dataset, each implemented as a dedicated DVC stage.

\begin{table}[t]
  \caption{The four variants the ESA dataset is prepared in.}
  \centering
  \small
  \begin{tabular}{lll}
    \toprule
    Positive class & Training split \\
    \midrule
    anomalies and rare events & positives removed \\
    anomalies and rare events & positives kept\\
    anomalies only            & positives removed \\
    anomalies only            & positives kept    \\
    \bottomrule
  \end{tabular}
  \label{tab:variants}
\end{table}

\section{Modelling}

Five distinct detectors are fitted on identical windowed input representations. 
Each model computes a Euclidean distance-based anomaly score, 
which is subsequently thresholded using the $3\sigma$ rule. 
Each detector is associated with a dedicated hyperparameter space.
Owing to compute and time constraints within the single-machine execution envelope, 
the parameter spaces for DENCAD, the DENCAST distance baseline, and the graph autoencoder are evaluated via random search, 
sampling 15 configurations per model. 
Hyperparameter search is conducted on the validation split using models trained exclusively on training data, 
optimized against the threshold-free Average Precision (AP) metric.

Following hyperparameter selection, 
each detector undergoes a final refitting step on the merged training and validation splits 
before being evaluated once on the isolated test set. 
Final comparisons are conducted using both threshold-free metrics: Average Precision (AP) and Area Under the Receiver Operating Characteristic Curve (ROC-AUC). 

Overall, autoencoder-based detectors slightly outperform distance-based baselines 
(K-Means and DENCAST distance) across nearly all dataset configurations under both metrics, 
whereas performance differences between DENCAST distance and K-Means remain marginal. 
Among the autoencoder-based architectures, no single model strictly dominates, 
and incorporating graph-level topology does not consistently justify its added complexity:
the plain autoencoder, which operates on isolated observations without graph structure, 
achieves the top performance on three of the six dataset variants under both metrics. 
Nevertheless, DENCAD secures the highest Average Precision on both ESA scenarios where rare events are treated as nominal,
demonstrating a performance gain over isolated baselines that is substantially more pronounced in AP than indicated by ROC-AUC.

\section{Evaluation}
\label{sec:eval}

\paragraph{Synthesis of empirical findings.}

All five evaluated models demonstrate suitability for the anomaly detection task, 
with autoencoder-based architectures exhibiting a slight overall performance advantage over distance-based baselines. 
However, the additional structural and computational complexity introduced by DENCAD does not 
yield a consistent or substantial performance gain to justify its overhead in these operational settings.

\section{Deployment}

The deliverable is not a serialised model but a pipeline that can rebuild one. It is a DVC
graph with one preprocessing and one splitting stage per dataset, one processing stage per
variant, and per variant a final fit for each detector preceded by its hyperparameter search. 
Every stage declares its code, data and parameter dependencies, so the lock file
records which inputs produced which report and a stale result cannot be mistaken for a
current one. All tunable quantities live in a single parameter file, so adding a variant is an entry in that file rather than a
change to the code. JSON reports are written as artefacts and the
dataset profiling lives in executed notebooks.

\end{document}
